% !TEX root = ../b-rg.tex
\chapter{Negation}\label{ch:negation}

Borlish has two clausal negators, and which one a clause takes depends on its finite verb. The auxiliary and copular verbs \bl{star} `be', \bl{er}, \bl{aïr} `have' and \bl{sceir} `there be' are negated with \bl{no}, placed before them. All other verbs are negated with \bl{noc}, placed after them. The two negators also behave differently with negative words such as \bl{alcun hom} `nobody' and \bl{jamay} `never': these replace \bl{noc}, but not \bl{no}. The future with \bl{volir} has its own negative verb, \bl{nolir} (Section~\ref{sec:vol-future}).

\section{Preverbal \bl{no}}\label{sec:neg-no}
The verbs \bl{star}, \bl{er}, \bl{aïr} and \bl{sceir} are negated by \bl{no} after the subject and before the verb. Any clitic pronoun that belongs to the verb itself comes between the two: \bl{ðe no lo son} `they are not (so)', \bl{la no l'aye} `she did not (have it)', \bl{Paravant no m'er nien tramis} `I wouldn't inherit anything before' (literally `before, nothing was being passed down to me'). Before a vowel it is elided to \bl{n'} (\bl{n'a}, \bl{n'ay}, \bl{n'au}, \bl{n'aun}, \bl{n'aye}, \bl{n'eu}). With the vowel-initial forms of `be', it contracts instead: \bl{no es} becomes \bl{no's}, and \bl{no er}, \bl{no ern} become \bl{no'r}, \bl{no'rn}. The first person singular \bl{no so} is also contracted to \bl{no's} in speech, so that \bl{no's} serves both persons; the full form is more careful.
    \glsen{Alcun hom n'a un cla.}{alˈkɪn ɔm na ɪn kla}{no person {\Neg=have} {\Indf} key}{Nobody has a key.}
    \glsen{No scey plu alcun laið.}{no xi pli alˈkɪn leθ}{{\Neg} {there.be} {no.more} no milk}{There's no milk left.}
Since the perfect, the progressive and the passive are formed with these verbs, \bl{no} negates all of them, standing before the auxiliary (\bl{Ðe n'aun nien ostendað tras dou meis} `They've posted nothing in two months'). Other adverbs come between the auxiliary and the participle.
    \glsen{Le no's ja ac yent.}{le nɔz ʒa ak jɛnt}{{\Third\Sg.\M} {\Neg=be} yet here {arrive.\Pst.\Ptcp}}{He hasn't arrived yet.}
    \glsen{Jo no's jamay eut a Rom.}{ʒo nɔz ʒaˈme awt a rɔm}{{\First\Sg} {\Neg=be} never {go.\Pst.\Ptcp} to Rome}{I've never been to Rome.}
The modal \bl{poðir} `can' may also take \bl{no} (\bl{Jo no poð lou affaçar} `I can't confront them'), but \bl{noc} after it is more common.

\section{Postverbal \bl{noc}}\label{sec:neg-noc}
Every other verb is negated by \bl{noc}, placed after the finite verb. Any object, complement or infinitive follows \bl{noc}. This includes the modal verbs, the synthetic future and the subjunctive.
    \glsen{Nos poðau noc y quarf trovar.}{nɔz poˈðo nɔk i kwarf troˈvar}{{\First\Pl} {can-\Ipfv.\First\Pl} {\Neg} {\Def} station {find-\Inf}}{We couldn't find the station.}
    \glsen{Jo crey noc ig le saif y vertað.}{ʒo kri nɔk aj le sef i vɛrˈtaθ}{{\First\Sg} believe {\Neg} {\Comp} {\Third\Sg.\M} {know.\Sbjv} {\Def} truth}{I don't think he knows the truth.}
Compare \bl{Nos farreu noc deunx cossy alcheðr} `We won't get so lucky again' and \bl{Mettau noc carrug avant naumagl} `Let's not get ahead of ourselves' (literally `Let's not put a cart before an ox'). The object clitics keep their place before the verb (\bl{Skig le keyottar; luy stou noc un douçur may} `Don't pamper him; he doesn't need more pudding').

In a question with inversion, \bl{noc} still follows the verb. A negative yes--no question is simply a negative statement with the final particle \bl{ou}: \bl{Tu la boðau noc ou?} `Didn't you call her?' (Chapter~\ref{ch:sentence-types}).

\subsection{Negative commands}
A negative command places \bl{noc} after the verb and after any pronoun that follows it (Section~\ref{sec:imperative-pronouns}): \bl{Deskieð tey noc} `Don't worry', \bl{Desmettað noc y sontiçtað\ldots} `Don't fire the backbone\ldots'. A full noun-phrase object follows \bl{noc}.
    \glsen{Toc noc l'oç!}{tɔk nɔk lɔts}{touch {\Neg} {\Def=\Dist.\Sg}}{Don't touch that!}
More informal, and more likely in modern speech, is \bl{skig} `avoid', the imperative of \bl{skighar}, with an infinitive. It is especially common with pronoun objects, which then stand before the infinitive in the ordinary way: \bl{Skig l'oç toccar!}, \bl{Skig le keyottar}.
% CHECK: a third-person clitic after a negative command (*Toc lo noc!) is
% still unattested; speakers seem to prefer skig + infinitive there.

\subsection{The negative future}
The future with \bl{volir} and an infinitive is not negated with \bl{noc}. Its negative is the verb \bl{nolir} `won't' with the infinitive, which also expresses refusal (\bl{Ðe noun nos aïðar} `They won't help us'; Section~\ref{sec:vol-future}). Its subjunctive \bl{nel} gives the negative of polite requests (\bl{Neleð y gran key roug appoyar} `Please don't press the big red button'; Section~\ref{sec:commands}).

\section{Negative words and negative concord}\label{sec:neg-words}
The negative words are \bl{alcun} `no, not any' with a noun, \bl{alcun hom} `nobody', \bl{nien} `nothing', \bl{necun} `neither, none', \bl{jamay} `never' and \bl{plu} `no more, not any longer'. \bl{Ja} `already, yet' means `yet' in a negative clause, and \bl{lon tan} `without also' is `not \ldots{} either' (\bl{e la poð noc viglar lon tan} `and she can't drive either').

With a verb that takes \bl{noc}, a negative word negates the clause on its own, and \bl{noc} is left out. This is so wherever the negative word stands: before the verb as subject, after it as object or postverbal subject, or as an adverb.
    \glsen{Alcun hom desir l'ou derran.}{alˈkɪn ɔm deˈzɪr lu deˈran}{no person want {\Def=egg} last}{Nobody wants the last egg.}
Compare \bl{Y molin produg alcun zucar\ldots} `The mill hasn't produced any sugar\ldots', \bl{Reman alcun laið} `There's no milk left' (literally `No milk remains'), and \bl{Y jammel se paroln jamay a lou meðes} `The twins never speak to each other'.

With the verbs that take \bl{no}, on the other hand, \bl{no} is kept alongside the negative word. The negatives do not cancel one another but reinforce each other, and any number of them may combine, as in \bl{No scey plu alcun laið} above.
    \glsen{La n'a lo dit a alcun hom.}{la na lo dɪt a alˈkɪn ɔm}{{\Third\Sg.\F} {\Neg=have} {\Third\Sg.\N.\Acc} {say.\Pst.\Ptcp} to no person}{She hasn't told anyone.}
Compare \bl{jo n'ay jamay lo vis} `I've never seen it', \bl{Ðe n'aun nien ostendað} `They've posted nothing' and, with the negative word as subject, \bl{Alcun hom n'au y lamp fornið con zygom} `No-one had installed a switch for the light' and \bl{Necun de lou no son prest} `Neither of them is ready'. So \bl{jamay} and the other negative words take the place of \bl{noc}, but never of \bl{no}.
% CHECK: is No reman alcun laið (no with a lexical verb) ever possible?

\section{Constituent negation}\label{sec:constituent-negation}
To negate a single constituent rather than the clause, \bl{noc} is placed before it: \bl{noc Theo} `not Theo', \bl{Noc y roug, y vert} `Not the red one, the green one', \bl{l'oç noc meyon amað} `the one I don't like' (Section~\ref{sec:object-relatives}). A quantified subject is negated in the same way. With a verb that takes \bl{noc}, the verb is then not negated a second time.
    \glsen{Noc toð gent accordn.}{nɔk tɔθ ʒɛnt aˈkɔrd.n̩}{{\Neg} all people {agree-\Third\Pl}}{Not everyone agrees.}
But where the verb is one that takes \bl{no}, including the auxiliary of the progressive, it is negated as well, by negative concord: \bl{Noc toð gent no son accordant} `Not everyone is agreeing'.
    \glsen{Noc toð gent no'rn cur dant.}{nɔk tɔθ ʒɛnt nɔrn kɪr dant}{{\Neg} all people {\Neg=\Cop.\Ipfv-\Third\Pl} attention {give-\Prs.\Ptcp}}{Not everyone was listening.}

\section{Negation and coordination}\label{sec:neg-coordination}
`Neither \ldots{} nor' is \bl{necun} \ldots{} \bl{ben}, where \bl{ben} is the ordinary word for `or'. After a verb that takes \bl{noc}, \bl{necun} replaces the negator. As a subject the pair stands before the verb, which is singular.
    \glsen{L'accatau necun pan ben kes.}{la.kaˈto neˈkɪn pan bɛn kɛz}{{\Third\Sg.\F=buy-\Pst} neither bread or cheese}{She didn't buy bread or cheese.}
    \glsen{Necun Ana ben Tom venoy.}{neˈkɪn ˈa.na bɛn tɔm veˈnɔj}{neither Ana or Tom {come-\Pst}}{Neither Ana nor Tom came.}
`And neither is X', after a negative clause, is \bl{ben} + X + \bl{necun}, with the verb left out.
    \glsen{Jo no's ig vant, ben Marc necun.}{ʒo nɔz aj vant bɛn mark neˈkɪn}{{\First\Sg} {\Neg=be} there {go-\Prs.\Ptcp} or Marc neither}{I'm not going, and neither is Marc.}
A negative verb phrase is left out after its subject, with only \bl{noc} remaining, or replaced by the pro-verb \bl{oc fair} `do so' with \bl{noc}: \bl{\ldots{} ma Theo oc fey noc}.
    \glsen{Marc e Halina venoirn, ma Theo noc.}{mark e ˌha.liˈna veˈnɔjrn̩ ma θeˈo nɔk}{Marc and Halina {come-\Pst.\Third\Pl} but Theo {\Neg}}{Marc and Halina came, but Theo didn't.}

\section{Negation in complex sentences}\label{sec:neg-complex}
With verbs of thinking and believing, the negation stays on the main verb, as in English \textit{I don't think\ldots}, and the complement clause then takes the subjunctive, usually introduced by \bl{ig} (\bl{Jo crey noc ig le saif y vertað}, above). A negated verb of knowing likewise takes the subjunctive, with \bl{ou} `whether' (\bl{Jo say noc ou l'ay figl} `I don't know whether he has children'; Section~\ref{sec:subjunctive}).

A negative purpose clause is introduced by \bl{neðan} `lest, so that \ldots{} not', with the subjunctive, or with an infinitive when its subject is the same as the main clause's (\bl{Gevou neðan cair!} `Careful not to fall!'; Section~\ref{sec:sbjv-adverbial}). `Without doing' is \bl{lon} with an infinitive (\bl{La parteu lon un parol digr} `She left without saying a word').

A negative question expecting a positive answer is contradicted not with \bl{ja} `yes' but with \bl{vis} `of course, indeed', often followed by a pro-verb with inversion: \bl{Tu la boðau noc ou? --- Vis oc fey jo!} `Didn't you call her? --- Yes, I did!' (Chapter~\ref{ch:sentence-types}).
